Furthermore, taking into account the models of Parolo et al.\cite{Par1} and Wang et al.\cite{Wan1}, we classify highly cited papers into the following five classes (FP/SB/Nm/G/O) using this index.
Among them, for the class Nm, we further establish a subclass No for subsequent analysis.
\begin{itemize}
\item FP: Flash in the pan: Sun et al. proposed $Gs=-0.600$ and $A^- = -0.300$ as the criteria for Flash in the pan. Referencing these values, we define papers with $Gs < -0.600$ and $A^- < -0.300$ as Flash in the pan.
\item SB: Sleeping beauty: Sun et al. did not provide criteria for Sleeping beauty using $A^+$, but in this study, we define papers with $Gs > 0.600$ and $A^+ > 0.300$ as Sleeping beauty as conditions contrasting with Flash in the pan.
\item Nm: Mature paper: For standard papers that have received citations to a certain extent and are not classified as Flash in the pan or Sleeping beauty, we set the conditions as $-0.600 \leq Gs \leq 0.600$, $A^- < 0$, $A^+ > 0$, and $|A^-| > |A^+|$, and classify papers meeting these as mature papers.
\item No: Mature papers with confirmed obsolescence: Among mature papers, we define a set of papers with $A^- < -0.1$ as those exhibiting a strong tendency towards obsolescence.
\item G: Growing paper: For papers in the growth stage of citations, we set the conditions as $-0.600 \leq Gs \leq 0.600$ and $A^- = 0$.
\item O: Others: Special citation shapes (such as revival papers) are considered to fall outside the above conditions.
\end{itemize}

